\clearpage
\begin{summarybox}
  \textbf{\textcolor{CSummary}{先看圆：整个结论的初中几何入口}}
  圆 $x^2+y^2=r^2$ 上，若外点 $P$ 的切线 $PA\perp PB$，则 $OA\perp PA$、$OB\perp PB$，故四边形 $OAPB$ 是矩形；再由 $OA=OB=r$，它是正方形。因此
  \[
    \boxed{OP^2=2r^2},\qquad
    \boxed{AB=OP=\sqrt2\,r},\qquad
    \boxed{AB\perp OP},
  \]
  且 $AB$ 与 $OP$ 互相平分。反过来，$OP=\sqrt2\,r$ 时，
  $PA^2=OP^2-r^2=r^2$，因此 $\angle APO=\angle OPB=45^\circ$，两条切线也垂直。所以轨迹确实是\textbf{整个圆} $\boxed{x^2+y^2=2r^2}$。

  \textbf{圆与椭圆的一眼对照：}
  \[
    \begin{aligned}
      \text{原曲线是圆 }x^2+y^2=r^2
         &\quad\Longrightarrow\quad OP^2=r^2+r^2=2r^2,\\
      \text{原曲线是椭圆 }\frac{x^2}{a^2}+\frac{y^2}{b^2}=1
         &\quad\Longrightarrow\quad OP^2=a^2+b^2.
    \end{aligned}
  \]
  \textbf{两个层次要分开：}圆的正方形直接用勾股定理；椭圆的平方和虽形式相似，却不能把正方形直接搬过去，必须另用椭圆相切条件证明。圆对应 $a=b=r$，恰是椭圆公式的特殊情况。

  \medskip

  \textbf{\textcolor{CSummary}{一句话核心}}
  椭圆 $\frac{x^2}{a^2}+\frac{y^2}{b^2}=1$（$a>b>0$）的\textbf{两条互相垂直的切线的交点} $P$，满足
  \[
    \boxed{OP^2=x_0^2+y_0^2=a^2+b^2},
  \]
  而且\textbf{反过来也成立}；轨迹是圆心 $O$、半径 $\sqrt{a^2+b^2}$ 的圆。

  \medskip
  \textbf{\textcolor{CSummary}{五种证明怎么选}}
  \begin{enumerate}[leftmargin=*, itemsep=0.25em]
    \item \textbf{门槛最低：}两条垂距 $\rightarrow$ 勾股定理 $\rightarrow$ 平方和不变。
    \item \textbf{速度最快：}斜率方程 $\rightarrow$ 韦达乘积 $-1$ $\rightarrow$ 圆方程。
    \item \textbf{步骤最少：}一般式 $Ax+By=C$ 与 $Bx-Ay=D$，利用相切条件两式相加。
    \item \textbf{常规通用：}设线 $\rightarrow$ 联立椭圆 $\rightarrow$ 判别式为 $0$ $\rightarrow$ 韦达 $\rightarrow$ 检查特殊情况与逆命题。
    \item \textbf{切点弦＋韦达（新增）：}切点弦 $AB$ $\rightarrow$ 只求 $x_1x_2,y_1y_2$ $\rightarrow$ 两条切线一般式的垂直条件；不求四组和与积。
  \end{enumerate}

  \medskip
  \textbf{\textcolor{CSummary}{新增方法：切点弦＋韦达的速用链}}
  \[
    \frac{x_0x}{a^2}+\frac{y_0y}{b^2}=1
    \quad\text{（切点弦）}.
  \]
  设 $A(x_1,y_1)$、$B(x_2,y_2)$，记 $D=b^2x_0^2+a^2y_0^2>0$，则
  \[
    \frac{x_1x_2}{a^4}=\frac{b^2-y_0^2}{D},\qquad
    \frac{y_1y_2}{b^4}=\frac{a^2-x_0^2}{D}.
  \]
  \[
    PA\perp PB
    \Longleftrightarrow\frac{x_1x_2}{a^4}+\frac{y_1y_2}{b^4}=0
    \Longleftrightarrow\boxed{x_0^2+y_0^2=a^2+b^2}.
  \]
  \textbf{记忆线索：}切点弦找切点，韦达只求\textbf{两个乘积}；两切线垂直看\textbf{一般式的系数}，不必求两条斜率。

  \medskip

  \medskip
  \textbf{\textcolor{CSummary}{补充发现一：为什么切点弦与切线代入后是同一个等式？}}
  设椭圆外点 $P(x_0,y_0)$ 引切线 $PA$、$PB$，切点分别是
  $A(x_1,y_1)$、$B(x_2,y_2)$。先区分\textbf{三条不同的直线}：
  \[
    \begin{aligned}
      AB:&\quad \frac{x_0x}{a^2}+\frac{y_0y}{b^2}=1
        &&\text{（切点弦，经过 }A,B\text{）},\\
      PA:&\quad \frac{x_1x}{a^2}+\frac{y_1y}{b^2}=1
        &&\text{（在 }A\text{ 处的切线，经过 }P\text{）},\\
      PB:&\quad \frac{x_2x}{a^2}+\frac{y_2y}{b^2}=1
        &&\text{（在 }B\text{ 处的切线，经过 }P\text{）}.
    \end{aligned}
  \]
  \textbf{第 1 组：}把 $A$ 代入切点弦 $AB$，与把 $P$ 代入切线 $PA$，分别得到
  \[
    \frac{x_0x_1}{a^2}+\frac{y_0y_1}{b^2}=1,
    \qquad
    \frac{x_1x_0}{a^2}+\frac{y_1y_0}{b^2}=1.
  \]
  \textbf{第 2 组：}把 $B$ 代入切点弦 $AB$，与把 $P$ 代入切线 $PB$，分别得到
  \[
    \frac{x_0x_2}{a^2}+\frac{y_0y_2}{b^2}=1,
    \qquad
    \frac{x_2x_0}{a^2}+\frac{y_2y_0}{b^2}=1.
  \]
  \textbf{为什么相同？}只是乘法交换律 $x_0x_i=x_ix_0$、$y_0y_i=y_iy_0$（$i=1,2$）。
  因此，\textbf{已知 $A$ 在椭圆上}时，$A$ 在 $P$ 对应的切点弦上，
  与 $P$ 在 $A$ 处切线上，是同一个坐标关系；$B$ 也一样。
  \textbf{注意：}代入后等式相同，\emph{不代表} $AB$、$PA$、$PB$ 是同一条直线。

  \medskip
  \textbf{\textcolor{CSummary}{补充发现二：两式相减定方向，两式相加找中点}}
  真正不同的两个等式为
  \[
    \frac{x_0x_1}{a^2}+\frac{y_0y_1}{b^2}=1,
    \qquad
    \frac{x_0x_2}{a^2}+\frac{y_0y_2}{b^2}=1.
  \]
  \textbf{相减}消去右边的 $1$，得到
  \[
    \frac{x_0(x_1-x_2)}{a^2}
      +\frac{y_0(y_1-y_2)}{b^2}=0.
  \]
  当 $y_0\ne0$、$x_1\ne x_2$ 时，移项并相除得到
  \[
    \boxed{k_{AB}=\frac{y_1-y_2}{x_1-x_2}
      =-\frac{b^2x_0}{a^2y_0}}.
  \]
  \textbf{几何意义：}相减得到的是\textbf{切点弦 $AB$ 的方向}，不是 $PA$ 或 $PB$ 的方向。
  若 $y_0=0$，切点弦为 $x=a^2/x_0$（竖直）；若 $x_0=0$，切点弦为
  $y=b^2/y_0$（水平），可以直接从切点弦方程判断，不必强行求斜率。

  \textbf{相加}得到
  \[
    \frac{x_0(x_1+x_2)}{a^2}
      +\frac{y_0(y_1+y_2)}{b^2}=2.
  \]
  设 $AB$ 的中点为 $M(x_m,y_m)$，其中
  $x_m=(x_1+x_2)/2$，$y_m=(y_1+y_2)/2$，则
  \[
    \boxed{\frac{x_0x_m}{a^2}+\frac{y_0y_m}{b^2}=1}.
  \]
  \textbf{几何意义：}中点 $M$ 仍在直线 $AB$ 上；这是已有两式的结果，
  \textbf{并没有提供新的独立条件}，单靠相加、相减还不能推出蒙日圆。

  \medskip
  \textbf{\textcolor{CSummary}{补充发现三：要证明蒙日圆，还缺哪个条件？}}
  前面的代入等式只反映\textbf{切点与切线的关系}；真正新增的条件是
  $PA\perp PB$。由两条切线的一般式，其垂直条件为
  \[
    \boxed{\frac{x_1x_2}{a^4}+\frac{y_1y_2}{b^4}=0}.
  \]
  把前面韦达定理的两个乘积代入，得到
  \[
    \frac{b^2-y_0^2}{D}+\frac{a^2-x_0^2}{D}=0
    \ \Longrightarrow\ 
    \boxed{x_0^2+y_0^2=a^2+b^2},
    \qquad D=b^2x_0^2+a^2y_0^2>0.
  \]
  \textbf{解题意识：}先看最终需要\emph{什么条件、哪些量}。
  此题新增的是“两切线垂直”，所需的是\textbf{两个坐标乘积}；
  同一关系反复代入不能代替新条件，也不必把四组和与积全部算出。

  \medskip
  \textbf{\textcolor{CSummary}{一眼核对：相同等式不等于同一条直线}}
  对椭圆 $\frac{x^2}{4}+y^2=1$，取 $P(2,1)$、$A(2,0)$、$B(0,1)$。
  此时
  \[
    AB:\ \frac{x}{2}+y=1,\qquad
    PA:\ x=2,\qquad PB:\ y=1.
  \]
  它们显然是三条不同的直线。把 $A$ 代入 $AB$ 与把 $P$ 代入 $PA$，
  都得到 $1=1$；把 $B$ 代入 $AB$ 与把 $P$ 代入 $PB$ 也一样。
  \textbf{几何意义：}同一等式说的是“某点在某条直线上”，不能据此把三条直线混为一谈。

  \textbf{\textcolor{CSummary}{考场必查}}
  \begin{itemize}[leftmargin=*, itemsep=0.25em]
    \item $a^2+b^2$ 是\textbf{平方和}，不要混成焦距关系的平方差。
    \item 使用韦达定理前先检查 $x_0^2\ne a^2$；竖直切线 $x=\pm a$ 不能写斜率。
    \item 要证明“轨迹是整个圆”，\textbf{正向和反向都要证明}。
    \item 切点弦 $AB$ 不等于 $PA,PB$；使用切点弦法时只求 $x_1x_2,y_1y_2$，不要习惯性把四组和与积全部算出。
    \item 圆上点都在椭圆外：$OP^2=a^2+b^2>a^2$，而椭圆上点到原点的距离平方不超过 $a^2$。
  \end{itemize}
\end{summarybox}
